%!TEX root = ../main.tex
% =============================================================
%  CAPÍTULO 5 — JUSTIFICACIÓN ECONÓMICA (Entrega 5 · 15 % · 15/9)
% =============================================================

\chapter{Justificación económica}\label{cap:economica}

\begin{guia}[Guía de la cátedra: qué se evalúa en este capítulo]
Es el capítulo de mayor peso individual (15\,\%). Se evalúa como una
evaluación de proyectos de inversión:
\begin{enumerate}
    \item \textbf{Análisis de mercado y competidores}: tamaño, segmentos,
          precios de referencia.
    \item \textbf{Modelo de negocio}: B2B, B2C, B2G o interno; niveles de
          precio y fuentes de ingreso.
    \item \textbf{Inversión inicial} desglosada por rol y tiempo, más
          infraestructura.
    \item \textbf{Costos operativos} recurrentes.
    \item \textbf{Escenarios} pesimista, base y optimista, con flujo de caja
          proyectado a 3 a 5 años.
    \item \textbf{Indicadores}: \gls{van}, \gls{tir} y período de repago,
          con una tasa de descuento construida y justificada con fuentes.
    \item \textbf{Supuestos y restricciones} explícitos.
\end{enumerate}
Todo número con su origen: cotización, tarifa pública, estadística
citada o estimación propia declarada como tal.
\end{guia}

\completar{Desarrollo del capítulo.}

\section{\completar{Nombre de la sección}}

\completar{Texto.}

\begin{table}[H]
    \centering
    \caption{Inversión inicial desglosada (ejemplo de estructura)}
    \label{tab:inversion}
    \small
    \begin{tabularx}{\textwidth}{L r r r L}
        \toprule
        \tb{Concepto} & \tb{Cantidad} & \tb{Unitario (USD)} & \tb{Total (USD)} & \tb{Fuente} \\
        \midrule
        Desarrollador full stack (meses) & 6 & \completar{2.500} & \completar{15.000} & \completar{encuesta salarial citada} \\
        Científico de datos (meses) & 4 & \completar{...} & \completar{...} & \completar{...} \\
        Diseño \gls{ux} (horas) & 80 & \completar{...} & \completar{...} & \completar{...} \\
        Infraestructura inicial & 1 & \completar{...} & \completar{...} & \completar{tarifa pública del proveedor} \\
        \midrule
        \multicolumn{3}{r}{\tb{Total}} & \tb{\completar{...}} & \\
        \bottomrule
    \end{tabularx}
\end{table}

\section{\completar{Nombre de la sección}}

\completar{Texto.}

\begin{table}[H]
    \centering
    \caption{Flujo de caja proyectado por escenario, en USD (ejemplo de estructura)}
    \label{tab:flujo-caja}
    \small
    \begin{tabularx}{\textwidth}{l R R R R R R}
        \toprule
        \tb{Escenario} & \tb{Año 0} & \tb{Año 1} & \tb{Año 2} & \tb{Año 3} & \tb{Año 4} & \tb{Año 5} \\
        \midrule
        Pesimista & $-$\completar{...} & \completar{...} & \completar{...} & \completar{...} & \completar{...} & \completar{...} \\
        Base      & $-$\completar{...} & \completar{...} & \completar{...} & \completar{...} & \completar{...} & \completar{...} \\
        Optimista & $-$\completar{...} & \completar{...} & \completar{...} & \completar{...} & \completar{...} & \completar{...} \\
        \bottomrule
    \end{tabularx}
\end{table}

\begin{ejemplo}[Ejemplo: construcción de la tasa de descuento]
La tasa de descuento \gls{sym:r} se construye sumando una tasa libre de
riesgo (bono del Tesoro de EE.\,UU.\ a 10 años), el riesgo país (índice
EMBI) y una prima por riesgo de proyecto en etapa temprana. Cada
componente se cita con su fuente y fecha de consulta.
\end{ejemplo}

El \gls{sym:van} se calcula como
\begin{equation}
    \mathit{VAN} = -I_0 + \sum_{t=1}^{n} \frac{\gls{sym:fc}}{(1 + r)^{t}},
    \label{eq:van}
\end{equation}
donde $I_0$ es la inversión inicial, $FC_t$ el flujo de caja neto del
año $t$ y $r$ la tasa de descuento. La \gls{tir} es la tasa que anula el
\gls{van}, y el período de repago es el primer año en que el flujo
acumulado se vuelve positivo.

\begin{table}[H]
    \centering
    \caption{Indicadores financieros por escenario (ejemplo de estructura)}
    \label{tab:indicadores-financieros}
    \small
    \begin{tabular}{l r r r}
        \toprule
        \tb{Escenario} & \tb{VAN (USD)} & \tb{TIR} & \tb{Repago} \\
        \midrule
        Pesimista & \completar{...} & \completar{...\,\%} & \completar{n} años \\
        Base      & \completar{...} & \completar{...\,\%} & \completar{n} años \\
        Optimista & \completar{...} & \completar{...\,\%} & \completar{n} años \\
        \bottomrule
    \end{tabular}
\end{table}
